% Consolidated 18-capacitance-energy from 309dbbd2c1c3ab17ea53128b49313f224ecea5a5; original TeX fragments preserved.

% === BEGIN 18-capacitance-energy-101-112.tex blob cf69b9dd7b76ad8ce162c156eaeeee97baa36b84 ===
% Irodov Part III §3.3 problems 101--112, user PDF pp.114--116,
% printed pp.105--107. Capacitor figure 3.16--3.17 inspected.
\electrocapproblem{101}{105}
半径为 $R_1$ 的孤立金属球外表紧贴同心介质壳，
介质相对介电常数为 $\varepsilon$，外半径 $R_2$；
其外部真空。求金属球相对于无穷远的电容。
\exsolution
设球总自由电荷为 $q$。高斯定律给
$D_r=q/(4\pi r^2)$，故
$R_1<r<R_2$ 时
$E_r=q/(4\pi\varepsilon_0\varepsilon r^2)$，
$r>R_2$ 时
$E_r=q/(4\pi\varepsilon_0r^2)$。
球电势
\[
\varphi(R_1)=\frac{q}{4\pi\varepsilon_0}
\left[\frac1\varepsilon
\left(\frac1{R_1}-\frac1{R_2}\right)
+\frac1{R_2}\right].
\]
由 $C=q/\varphi(R_1)$：
\[
\boxed{C=
\frac{4\pi\varepsilon_0\varepsilon R_1}
{1+(\varepsilon-1)R_1/R_2}}.
\]
\excheck $R_2=R_1$ 时退化为真空导体球电容
$4\pi\varepsilon_0R_1$；
$R_2\to\infty$ 时得到
$4\pi\varepsilon_0\varepsilon R_1$。

\electrocapproblem{102}{105}
两个空气平板电容器电容均为 $C$，
串联接在电动势 $\mathcal E$ 不变的电池两端。
随后将其中一个电容器完全填充
相对介电常数为 $\varepsilon$ 的介质。
问该电容器的电场减小多少倍，
并求经电池流过的电荷量。
\exsolution
原串联等效电容 $C_{\rm eq,0}=C/2$，
串联电荷 $Q_0=C\mathcal E/2$，
每个电容器各承压 $\mathcal E/2$。
填充后该电容器电容为 $\varepsilon C$，
新串联电容为
$C_{\rm eq}=\varepsilon C/(1+\varepsilon)$。
故它的新电压
$V_{\rm diel}=Q/(\varepsilon C)
=\mathcal E/(1+\varepsilon)$。
同一板距下 $E\propto V$，因此
\[
\boxed{\frac{E_{\rm new}}{E_{\rm old}}
=\frac2{\varepsilon+1}},
\quad
\text{原电场减小 }(\varepsilon+1)/2\text{ 倍}.
\]
通过电池的电荷量等于两次串联电荷差：
\[
\boxed{\Delta Q=
C\mathcal E\left[\frac{\varepsilon}{1+\varepsilon}
-\frac12\right]
=\frac{C\mathcal E(\varepsilon-1)}
{2(\varepsilon+1)}}.
\]
\excheck 电源连接未断开，必须保持
$\mathcal E$ 不变，而不能错误地保持串联电荷不变。

\electrocapproblem{103}{105--106}
平板电容器板面积为 $S$，间距由相邻
两层介质填满；厚度 $d_1,d_2$，
相对介电常数 $\varepsilon_1,\varepsilon_2$。
求电容及两层界面处束缚面电荷密度
（场从1层指向2层，外加电压 $V$）。
\exsolution
无自由界面电荷，法向位移量恒定 $D$，
两层 $E_i=D/(\varepsilon_0\varepsilon_i)$。
电压
$V=D[d_1/\varepsilon_1+d_2/\varepsilon_2]/\varepsilon_0$。
因此
\[
\boxed{C=
\frac{\varepsilon_0S}{d_1/\varepsilon_1+d_2/\varepsilon_2}}.
\]
两层极化分别为
$P_i=D(1-1/\varepsilon_i)$。
在界面上定义沿1层指向2层为正法向，
来自介质1的束缚面电荷为 $+P_1$，
来自介质2的为 $-P_2$，
相加得
\[
\sigma_b=P_1-P_2
=D\left(\frac1{\varepsilon_2}-\frac1{\varepsilon_1}\right).
\]
代入
$D=\varepsilon_0V/(d_1/\varepsilon_1+d_2/\varepsilon_2)$：
\[
\boxed{\sigma_b=
\frac{\varepsilon_0V(\varepsilon_1-\varepsilon_2)}
{\varepsilon_1d_2+\varepsilon_2d_1}}.
\]
\excheck 这是\emph{同一界面两侧束缚电荷之和}，
不是单取其中一层的表面极化电荷。

\electrocapproblem{104}{106}
平板电容器面积 $S$，板距 $d$，
介质相对介电常数沿电场方向线性地
由 $\varepsilon_1$ 增至 $\varepsilon_2>\varepsilon_1$。
求电容及带电 $q$ 时束缚体电荷密度
关于局部 $\varepsilon$ 的关系。
\exsolution
没有自由体电荷，$D=q/S$ 恒定，
$E(x)=D/[\varepsilon_0\varepsilon(x)]$。
令
$\varepsilon(x)=\varepsilon_1+
(\varepsilon_2-\varepsilon_1)x/d$，
则
\[
V=\frac{q}{\varepsilon_0S}
\int_0^d\frac{\mathrm dx}{\varepsilon(x)}
=\frac{qd}{\varepsilon_0S(\varepsilon_2-\varepsilon_1)}
\ln\frac{\varepsilon_2}{\varepsilon_1}.
\]
因此
\[
\boxed{C=
\frac{\varepsilon_0S(\varepsilon_2-\varepsilon_1)}
{d\ln(\varepsilon_2/\varepsilon_1)}}.
\]
$P=D(1-1/\varepsilon)$，
$\rho_b=-\mathrm dP/\mathrm dx$，故
\[
\boxed{\rho_b(x)=
-\frac{q}{S}\frac{\varepsilon_2-\varepsilon_1}
{d\,\varepsilon(x)^2}}.
\]
\excheck 若 $\varepsilon_2\to\varepsilon_1=\varepsilon$，
洛必达法则给 $C\to\varepsilon\varepsilon_0S/d$、
$\rho_b\to0$。

\electrocapproblem{105}{106}
半径为 $R_1,R_2$ 的同心球形电容器，
$R_2>R_1$。填充介质相对介电常数随半径
变化为 $\varepsilon(r)=a/r$，
其中 $a$ 为具有长度量纲的常数。
求电容。
\exsolution
径向位移
$D_r=q/(4\pi r^2)$，
电场
$E_r=D_r/[\varepsilon_0(a/r)]
=q/(4\pi\varepsilon_0 a r)$。
电压
\[
V=\int_{R_1}^{R_2}E_r\,\mathrm dr
=\frac{q}{4\pi\varepsilon_0a}
\ln\frac{R_2}{R_1}.
\]
因此
\[
\boxed{C=\frac{4\pi\varepsilon_0a}
{\ln(R_2/R_1)}}.
\]
\excheck 因 $\varepsilon=a/r$ 为无量纲，
$a$ 必须具有长度量纲。

\electrocapproblem{106}{106}
同轴圆柱电容器两层介质分别占
$R_1<r<R_2$ 与 $R_2<r<R_3$，
介电常数为 $\varepsilon_1,\varepsilon_2$，
击穿场强为 $E_{1m},E_{2m}$。
求增大电压时两层同时达到各自击穿场的条件。
\exsolution
每单位轴向长度自由电荷量记为 $\lambda$。
由柱面高斯定律
$D_r=\lambda/(2\pi r)$，
因此第 $i$ 层
$E_i(r)=|\lambda|/(2\pi\varepsilon_0\varepsilon_i r)$。
场强随半径增大而减小，
所以第1层最大场在 $R_1$，
第2层最大场在 $R_2$。
若两层同步击穿：
\[
\frac{|\lambda|}{2\pi\varepsilon_0\varepsilon_1R_1}
=E_{1m},\quad
\frac{|\lambda|}{2\pi\varepsilon_0\varepsilon_2R_2}
=E_{2m}.
\]
消去 $\lambda$：
\[
\boxed{\varepsilon_1R_1E_{1m}
=\varepsilon_2R_2E_{2m}}.
\]
\excheck 不能用相同的径向位置比较不同层击穿，
每层的危险位置是其各自内表面。

\electrocapproblem{107}{106}
图3.16的双层同轴圆柱电容器，
半径依次为 $R_1<R_2<R_3$，
介质参数同上。
若
$\varepsilon_1R_1E_{1m}<
\varepsilon_2R_2E_{2m}$，
求整体击穿电压。
\exsolution
由于第1层允许的线电荷阈值更小，
先在 $r=R_1$ 达到 $E_{1m}$，
此时
$|\lambda|=2\pi\varepsilon_0\varepsilon_1R_1E_{1m}$。
电压由两层径向积分：
\[
V=\frac{|\lambda|}{2\pi\varepsilon_0}
\left[\frac1{\varepsilon_1}\ln\frac{R_2}{R_1}
+\frac1{\varepsilon_2}\ln\frac{R_3}{R_2}\right].
\]
代入阈值：
\[
\boxed{V_{\rm breakdown}=
R_1E_{1m}
\left[\ln\frac{R_2}{R_1}
+\frac{\varepsilon_1}{\varepsilon_2}\ln\frac{R_3}{R_2}\right]}.
\]
\excheck 条件为严格小于，因此先击穿第1层；
同号等号时两层同时击穿。

\electrocapproblem{108}{106}
空气中两根等长平行细金属导线，
截面半径均为 $a$，轴间距 $b\gg a$。
求两导线\emph{每单位长度的互电容}。
\exsolution
设两线单位长度带自由电荷
$+\lambda,-\lambda$。
二维线电荷势
$\varphi\propto\ln r$。
以 $a\ll b$ 的近似在正线表面取
距正线 $a$、负线 $b$，
同理负线反向，得到电压
\[
V\simeq\frac{\lambda}{\pi\varepsilon_0}\ln\frac ba.
\]
故单位长度电容
\[
\boxed{C'=\frac{\lambda}{V}
\simeq\frac{\pi\varepsilon_0}{\ln(b/a)}}.
\]
\excheck 更精确的两圆柱公式
$C'=\pi\varepsilon_0/
\operatorname{arcosh}(b/(2a))$，
在 $b\gg a$ 化成上述结果。

\electrocapproblem{109}{106}
一根半径 $a$ 的长金属导线，
平行于无限金属导体平面，
轴线距平面 $b\gg a$。
求两者的每单位长度互电容。
\exsolution
用导体平面镜像法，
真实导线带线电荷 $+\lambda$，
其镜像为距真线 $2b$ 的 $-\lambda$。
真线与镜像的电势差约为
$(\lambda/\pi\varepsilon_0)\ln(2b/a)$，
但平面为对称零电势位置，
真线到平面的电压为其中一半：
\[
V\simeq\frac{\lambda}{2\pi\varepsilon_0}\ln\frac{2b}{a}.
\]
故
\[
\boxed{C'=\frac{\lambda}{V}
\simeq\frac{2\pi\varepsilon_0}{\ln(2b/a)}}.
\]
\excheck 对无限导体平面并未采用孤立双线的
完整电压，否则会误少一倍。

\electrocapproblem{110}{106}
两个半径同为 $a$ 的导体球，
球心距 $b\gg a$，
均浸于相对介电常数为 $\varepsilon$ 的均匀介质中。
求双球系统的互电容。
\exsolution
设两球带电 $+q,-q$，
在 $b\gg a$ 下近似各自球面电荷均匀。
正球电势约
$\varphi_+=kq/(\varepsilon a)-kq/(\varepsilon b)$，
负球电势相反。
因此
\[
V=\varphi_+-\varphi_-=
\frac{q}{2\pi\varepsilon_0\varepsilon}
\left(\frac1a-\frac1b\right).
\]
电容
\[
C=\frac qV=
\frac{2\pi\varepsilon_0\varepsilon}
{1/a-1/b}
\simeq\boxed{2\pi\varepsilon_0\varepsilon a}
\quad(b\gg a).
\]
\excheck $b$ 有限时球面电荷并非严格均匀；
上述 $1/b$ 项仅用于说明首阶近似。

\electrocapproblem{111}{106}
半径为 $a$ 的金属球，
球心距无限导体平面为 $\ell\gg a$。
求球—平面的互电容。
\exsolution
镜像法使金属球附近的电势
近似为球自身带电 $q$ 的电势
减去镜像 $-q$ 在距 $2\ell$ 处的电势：
\[
V\simeq\frac{q}{4\pi\varepsilon_0}
\left(\frac1a-\frac1{2\ell}\right).
\]
因此
\[
C\simeq\frac{4\pi\varepsilon_0}
{1/a-1/(2\ell)}
\longrightarrow
\boxed{4\pi\varepsilon_0a}\quad(\ell/a\to\infty).
\]
\excheck 球面电荷的畸变为更高阶修正，
本题只要求 $\ell\gg a$ 的主阶。

\electrocapproblem{112}{106--107}
原书图3.17(a)、(b)为两种电容器连接网络，
求 A、B 端等效电容，按图中标号
(a) $C_1,C_2,C_3$，(b) 每只电容均为 $C$。
\exsolution
(a) 图3.17(a) 中三只电容器各连接同一对
等势节点 A、B；它们两端的电势差相等，
故各自电荷 $Q_i=C_iV$，
总电荷 $Q=\sum_iQ_i$，
因此
\[
\boxed{C_{AB}^{(a)}=C_1+C_2+C_3}.
\]
(b) 图3.17(b) 的四个边电容构成
对称桥式结构；沿对称中心线的节点电势相等，
可将桥的中间电容视为无压差。
剩余两条支路各有两个相等电容串联，
每支电容 $C/2$，两支并联：
\[
\boxed{C_{AB}^{(b)}
=\frac C2+\frac C2=C}.
\]
\excheck 基于原书图3.17网络连线，
不是将四只电容全串联或全并联；
按 $C_i$ 同等电容极限可复核。
% === END 18-capacitance-energy-101-112.tex ===

% === BEGIN 18-capacitance-energy-113-121.tex blob 5808fb25bcc201507ed321b39eae11738067cb6c ===
% Irodov Part III §3.3 113--121; source Fig 3.18--3.23 inspected,
% user PDF pp.116--117 (printed pp.107--108).
\electrocapproblem{113}{107}
图3.18中四块面积相同为 $S$ 的平行导体板，
自上至下相邻间距都是 $d$。按图(a)、(b)
所示导线连接后，求 A、B 间电容。
\exsolution
记每个相邻间隙的平板电容
$C_0=\varepsilon_0S/d$。
(a) 图(a)把第1、3板通过左端导线短接，
第2板接 A，第4板接 B。
第1—2和第2—3间两个 $C_0$
并联成 $2C_0$，它再与第3—4间的
$C_0$ 串联，故
\[
\boxed{C_a=
\frac{(2C_0)C_0}{2C_0+C_0}
=\frac{2\varepsilon_0S}{3d}}.
\]
(b) 图(b)把第1、4板通过左端导线短接，
第2板接 A，第3板接 B。
第1—2间 $C_0$ 和第3—4间 $C_0$
以中间浮置导体节点形成串联支路，
等效 $C_0/2$；同时第2—3间
$C_0$ 直接跨 A、B 并联，故
\[
\boxed{C_b=C_0+\frac{C_0}{2}
=\frac{3\varepsilon_0S}{2d}}.
\]
\excheck 本题需要先按图识别短接的
\emph{同势金属板}，才能进行串并联化简；
不能把四块板按几何顺序简单串联。

\electrocapproblem{114}{107}
两电容器 $C_1=1.0\,\mathrm{\mu F}$，
最大耐压 $V_1=6.0\,\mathrm{kV}$；
$C_2=2.0\,\mathrm{\mu F}$，
最大耐压 $V_2=4.0\,\mathrm{kV}$。
串联后的最大安全总电压是多少？
\exsolution
串联时两电容器携带相同电荷 $Q$，
各自电压 $U_i=Q/C_i$。
安全条件为
$Q\le C_1V_1$ 与 $Q\le C_2V_2$，
所以
\[
Q_{\max}=\min(1.0\times6.0,\,
2.0\times4.0)\,\mathrm{mC}
=6.0\,\mathrm{mC}.
\]
总电压为
\[
V_{\max}=Q_{\max}
\left(\frac1{C_1}+\frac1{C_2}\right)
=V_1\left(1+\frac{C_1}{C_2}\right).
\]
因此
\[
\boxed{V_{\max}=9.0\,\mathrm{kV}}.
\]
\excheck 先到耐压极限的是 $C_1$；
此时 $C_2$ 电压为3 kV，低于4 kV。

\electrocapproblem{115}{107}
原书图3.19：一个电动势
$\mathcal E=110\,\mathrm V$ 的电源，
上支路两个相同电容 $C_1$ 串联，
各串联节点到下导线 B 另接电容 $C_2$；
末端节点为 A。
已知 $C_2/C_1=\eta=2.0$。
求 $\varphi_A-\varphi_B$。
\exsolution
取下端 B 电势为零，
电源上端电势 $\mathcal E$。
把中间节点电势记为 $x$，
末端 A 电势为 $y$。
未由外部注入自由电荷的浮置节点需净电荷为零：
\[
C_1(x-\mathcal E)+C_1(x-y)+C_2x=0,
\]
\[
C_1(y-x)+C_2y=0.
\]
用 $\eta=C_2/C_1$ 化简：
$(2+\eta)x-y=\mathcal E$、
$x=(1+\eta)y$。
相消得到
\[
\boxed{y=\varphi_A-\varphi_B=
\frac{\mathcal E}{1+3\eta+\eta^2}}.
\]
代入 $\eta=2$：
\[
\boxed{\varphi_A-\varphi_B=
\frac{110}{11}\,\mathrm V=10\,\mathrm V}.
\]
\excheck 若误把两个 $C_2$ 直接并联，
会破坏中间节点电荷守恒，
无法得到该分母。

\electrocapproblem{116}{107}
图3.20：一个从左向右无限重复的电容梯形网络，
每节沿上支路有电容 $C$，
节点到公共下支路又有电容 $C$。
求左端与下公共导线之间的总等效电容。
\exsolution
设整个无穷网络的等效电容为 $X$。
第一只上支路电容 $C$ 后，
从新节点向下有一个电容 $C$，
并联上与原网络形状完全相同的后续部分，
其电容亦为 $X$，故节点到地等效为 $C+X$。
因此
\[
\frac1X=\frac1C+\frac1{C+X}.
\]
令 $x=X/C$，整理：
$x=(1+x)/(2+x)$，
即
$x^2+x-1=0$。
正的物理解为
\[
\boxed{X=\frac{\sqrt5-1}{2}C\simeq0.618C}.
\]
\excheck 必须取正根；
另一个根 $-(1+\sqrt5)/2$ 不能为被动电容的等效值。

\electrocapproblem{117}{107--108}
图3.21的 A、B 间依次串接
电容器 $C_1=1.0\,\mathrm{\mu F}$、
电源 $\mathcal E=10\,\mathrm V$、
电容器 $C_2=2.0\,\mathrm{\mu F}$。
已知 $\varphi_A-\varphi_B=5.0\,\mathrm V$。
求两只电容器各自的电压大小。
\exsolution
两只电容器串联而中间电源可施加电势跃变，
两只电容器的电荷模相同，设为 $Q$。
按原书图3.21的电源极性，
两电容器两端电压和
$U_1+U_2=(\varphi_A-\varphi_B)+\mathcal E$。
因此
\[
Q\left(\frac1{C_1}+\frac1{C_2}\right)
=5.0+10.0=15.0\,\mathrm V,
\]
\[
Q=\frac{(15.0)C_1C_2}{C_1+C_2}
=10\,\mathrm{\mu C}.
\]
所以
\[
\boxed{U_1=\frac Q{C_1}=10\,\mathrm V},
\qquad
\boxed{U_2=\frac Q{C_2}=5.0\,\mathrm V}.
\]
\excheck 带电状态与电源极性相关；
电压\emph{模}为正，代数电压沿回路应保留方向。

\electrocapproblem{118}{108}
图3.22：一个闭合环路中两只电容器
$C_1,C_2$ 与两个电动势
$\mathcal E_1,\mathcal E_2$ 按图中方向连接。
求每只电容器左极板相对于右极板的电势差。
\exsolution
设电容器极板代数电荷模为 $Q$，
在无旁路充电的串联闭环中各自绝对值相同。
沿图3.22所取的电动势方向，
两个电源的代数贡献为
$\mathcal E_2-\mathcal E_1$，
所以
\[
Q\left(\frac1{C_1}+\frac1{C_2}\right)
=\mathcal E_2-\mathcal E_1.
\]
对两个电容器左至右的电压，
依图中板极性设定：
\[
\boxed{V_1=\frac{\mathcal E_2-\mathcal E_1}
{1+C_1/C_2}},\qquad
\boxed{V_2=\frac{\mathcal E_1-\mathcal E_2}
{1+C_2/C_1}}.
\]
\excheck 两只电容器的左—右电压并非
同时同号；若交换源的正负端，
上述两个式子的代数符号须相应改变。

\electrocapproblem{119}{108}
仍在图3.22的电源与电容串联闭环中，
求每只电容器上的电荷\emph{模}。
\exsolution
由图中两个电源的代数电动势差
$\Delta\mathcal E=\mathcal E_2-\mathcal E_1$，
稳态的电容回路平衡为
$Q/C_1+Q/C_2=|\Delta\mathcal E|$。
于是
\[
\boxed{Q_1=Q_2=
\frac{C_1C_2}{C_1+C_2}
|\mathcal E_2-\mathcal E_1|}.
\]
\excheck 两只电容器带电量大小相等，
且 $\mathcal E_1=\mathcal E_2$ 时为零；
这是在未预充电、无其它净节点电荷条件下的结果。

\electrocapproblem{120}{108}
图3.23的四臂电容桥，
各臂电容为 $C_1,C_2,C_3,C_4$，
上下接电动势 $\mathcal E$，
取两个中间节点为 A、B。
求 $\varphi_A-\varphi_B$，以及零电势差条件。
\exsolution
电容桥两条并联支路分别为
$C_1$ 与 $C_3$ 串联、
$C_2$ 与 $C_4$ 串联；
选图中上下源极性，
在支路的电压分配公式中
\[
\varphi_A=\frac{C_2}{C_1+C_2}\mathcal E
\quad\text{与}\quad
\varphi_B=\frac{C_4}{C_3+C_4}\mathcal E
\]
在原书节点标号约定下，
对应两段电势差整理为
\[
\boxed{\varphi_A-\varphi_B=
\mathcal E\frac{C_2C_3-C_1C_4}
{(C_1+C_2)(C_3+C_4)}}.
\]
桥平衡条件是分子为零：
\[
\boxed{C_1/C_2=C_3/C_4}.
\]
\excheck 图中四个电容的编号不是
顺序围绕一个矩形而不变；
应按每一串联支路中
浮动节点的电荷守恒确定电压。

\electrocapproblem{121}{108}
一只已充至 $V=110\,\mathrm V$ 的
电容器 $C_1=1.0\,\mathrm{\mu F}$，
并联接到另外两只原先不带电、串联的
电容器 $C_2=2.0\,\mathrm{\mu F}$、
$C_3=3.0\,\mathrm{\mu F}$ 的两端。
求通过连接导线的电荷量。
\exsolution
设两只原先未带电的串联电容器最终带电荷模
$q$，其所分配电压为
$V_f=q(1/C_2+1/C_3)$。
$C_1$ 最终带电量为 $C_1V_f$，
而整个外部节点总电荷守恒：
$C_1V=C_1V_f+q$。
联立可得
\[
C_1V=q\left[1+C_1
\left(\frac1{C_2}+\frac1{C_3}\right)\right],
\]
即
\[
\boxed{q=\frac{V}
{1/C_1+1/C_2+1/C_3}}.
\]
代入数值：
\[
\boxed{q=\frac{110}{1+1/2+1/3}\,\mathrm{\mu C}
=60\,\mathrm{\mu C}=0.060\,\mathrm{mC}}.
\]
\excheck 电容值以 $\mu$F 计时，
式中 $V/(\mu\mathrm F^{-1})$ 为 $\mu$C；
通过任一串联连接导线的累计电荷量相同。
% === END 18-capacitance-energy-113-121.tex ===

% === BEGIN 18-capacitance-energy-122-130.tex blob 4f2af64594b5935c9b1ae26c09f9b17b0fc30101 ===
% Irodov Part III §3.3 122--130; PDF pp.117--119,
% printed pp.108--110. Fig.3.24--3.29 checked.
\electrocapproblem{122}{108}
图3.24：左、右两个电源电动势都为
$\mathcal E$，电容器 $C_1$ 接在右上分支，
$C_2$ 接在中间竖支路。原先开关 Sw 断开，
闭合后求沿图中箭头方向流过截面1、2的电荷量。
\exsolution
取底部左侧公共点为零电势。
闭合前，左侧电源未接入，
$C_1,C_2$ 串联接右侧电源，
初始串联电荷
$Q_0=\mathcal E C_1C_2/(C_1+C_2)$。
闭合后左电源将中间上节点钳在
$+\mathcal E$，右电源上端同样为
$+\mathcal E$，故 $C_1$ 两端电压归零，
$C_2$ 两端电压为 $\mathcal E$。
截面1按左箭头的累计电荷
等于 $C_2$ 的最终电荷：
\[
\boxed{q_1=\mathcal E C_2}.
\]
截面2按图中右箭头，$C_1$ 的
初有电荷 $Q_0$ 被消去，对该方向贡献为负：
\[
\boxed{q_2=-\frac{\mathcal E C_1C_2}{C_1+C_2}}.
\]
\excheck 截面电荷是开关动作前后的\emph{累计通过量}，
不等于最后流动的稳态电流（稳态电流为零）。

\electrocapproblem{123}{108}
图3.25：左侧两电源均为
$\mathcal E=60\,\mathrm V$，右侧两只电容器
$C_1=2.0\,\mathrm{\mu F}$、
$C_2=3.0\,\mathrm{\mu F}$。
闭合中间 Sw 后，求沿箭头1、2、3
的累计电荷量。
\exsolution
闭合前，右侧串联电容器跨越两个电源
给出的总电压 $2\mathcal E$，
故每只电容器的电荷模
$Q_0=2\mathcal E C_1C_2/(C_1+C_2)$。
闭合后两只电容器各跨一只电动势
$\mathcal E$，分别带电
$Q_{1f}=\mathcal E C_1$、
$Q_{2f}=\mathcal E C_2$。
由各导线截面的电荷连续性，依原书箭头：
\[
\boxed{q_1=\mathcal E C_1-Q_0
=\mathcal E C_1\frac{C_1-C_2}{C_1+C_2}},
\]
\[
\boxed{q_2=-\mathcal E C_2
\frac{C_2-C_1}{C_1+C_2}
=\mathcal E C_2\frac{C_1-C_2}{C_1+C_2}},
\]
\[
\boxed{q_3=\mathcal E(C_2-C_1)}.
\]
代入数据：
\[
\boxed{q_1=-24\,\mathrm{\mu C},\qquad
q_2=-36\,\mathrm{\mu C},\qquad
q_3=+60\,\mathrm{\mu C}}.
\]
\excheck 两个负号说明实际电荷流过截面1、2时
与图中所指箭头相反；第3截面为正。

\electrocapproblem{124}{108--109}
原书图3.26：上支路为
$C_1,C_2$ 两只串联电容器，
中点 A 通过电容器 $C_3$ 连接下方 B，
底部分别有电动势 $\mathcal E_1,\mathcal E_2$。
求 $\varphi_A-\varphi_B$。
\exsolution
取 $\varphi_B=0$。
按图中电源极性，
$C_1$ 左端电势为 $-\mathcal E_1$，
$C_2$ 右端电势为 $+\mathcal E_2$。
A 是孤立浮置金属节点，初始无净自由电荷，
其三个电容器极板电荷代数和为零：
\[
C_1(\varphi_A+\mathcal E_1)
+C_2(\varphi_A-\mathcal E_2)
+C_3\varphi_A=0.
\]
于是
\[
\boxed{\varphi_A-\varphi_B=
\frac{C_2\mathcal E_2-C_1\mathcal E_1}
{C_1+C_2+C_3}}.
\]
\excheck $C_3$ 越大，中间节点相对于 B
的电势差越小，符合把 A 拉向 B 电势的直觉。

\electrocapproblem{125}{109}
原书图3.27：三个电源
$\mathcal E_1,\mathcal E_2,\mathcal E_3$
沿一竖支路串接，三只电容器
$C_1,C_2,C_3$ 分别把各层节点
1、2、3 接到公共点 O。
规定 $\varphi_O=0$，求
$\varphi_1,\varphi_2,\varphi_3$。
\exsolution
按图中各电动势极性设定，
三个左侧节点之间满足
$\varphi_i=\Phi-\mathcal E_i$，
其中 $\Phi$ 为共同参考偏移。
整个左侧电源组与外界通过电容耦合，
总净电荷守恒：
\[
C_1\varphi_1+C_2\varphi_2+C_3\varphi_3=0.
\]
故
$\Phi=(C_1\mathcal E_1+C_2\mathcal E_2+
C_3\mathcal E_3)/(C_1+C_2+C_3)$。
定义 $C_\Sigma=C_1+C_2+C_3$，得到
\[
\boxed{\varphi_1=
\frac{C_2\mathcal E_2+C_3\mathcal E_3-
(C_2+C_3)\mathcal E_1}{C_\Sigma}},
\]
\[
\boxed{\varphi_2=
\frac{C_1\mathcal E_1+C_3\mathcal E_3-
(C_1+C_3)\mathcal E_2}{C_\Sigma}},
\]
\[
\boxed{\varphi_3=
\frac{C_1\mathcal E_1+C_2\mathcal E_2-
(C_1+C_2)\mathcal E_3}{C_\Sigma}}.
\]
\excheck 将三式分别乘 $C_1,C_2,C_3$ 再相加，
应得零，满足浮置导体组的总电荷守恒。

\electrocapproblem{126}{109}
原书图3.28电容桥：
A—B 之间上支路依次为 $C_1,C_2$，
下支路依次为 $C_2,C_1$；
两串联支路中点通过一只 $C_3$ 桥接。
求总等效电容。
\exsolution
设 A 电势为 $V$，B 为零。
上下支路的中点电势分别为 $x,y$。
电荷守恒：
\[
C_1(x-V)+C_2x+C_3(x-y)=0,
\]
\[
C_2(y-V)+C_1y+C_3(y-x)=0.
\]
两式相加得 $x+y=V$；
相减得到
\[
x-y=\frac{(C_1-C_2)V}{C_1+C_2+2C_3}.
\]
A 端净电荷为
$Q=C_1(V-x)+C_2(V-y)$，
代入得
\[
C_{AB}=\frac QV
=\frac{2C_1C_2+C_3(C_1+C_2)}
{C_1+C_2+2C_3}.
\]
因此
\[
\boxed{C_{AB}=
\frac{2C_1C_2+C_3(C_1+C_2)}
{C_1+C_2+2C_3}}.
\]
\excheck $C_3\to0$ 时恢复
两条串联支路并联
$2C_1C_2/(C_1+C_2)$；
$C_1=C_2$ 时桥中点同势，结果为 $C_1$，
不依赖 $C_3$。

\electrocapproblem{127}{109}
图3.29(a)、(b)、(c)：
边长 $a$ 的正方形四顶点分别放置
大小均为 $q$ 的点电荷；(a) 全为正，
(b) 相邻顶点正负交错，
(c) 上两顶点正、下两顶点负。
求每种构型四电荷的相互作用总能。
\exsolution
离散点电荷的相互作用能
$W=k\sum_{i<j}q_iq_j/r_{ij}$，
正方形有4条边（距离 $a$）
及2条对角线（距离 $\sqrt2a$）。
(a) 全部相同，4条边贡献
$4kq^2/a$、对角线贡献
$\sqrt2kq^2/a$，所以
\[
\boxed{W_a=(4+\sqrt2)\frac{q^2}
{4\pi\varepsilon_0a}}.
\]
(b) 四条边均为异号，贡献
$-4kq^2/a$；两条对角线同号，
贡献 $+\sqrt2kq^2/a$：
\[
\boxed{W_b=(\sqrt2-4)\frac{q^2}
{4\pi\varepsilon_0a}}.
\]
(c) 两条水平边同号、两条竖直边异号，
边贡献抵消；两条对角线异号：
\[
\boxed{W_c=-\sqrt2\frac{q^2}
{4\pi\varepsilon_0a}}.
\]
\excheck 相互作用能只对每一\emph{无序电荷对}
计数一次，不加点电荷本身的无穷自能。

\electrocapproblem{128}{109}
无限直线链由 $q,-q,q,-q,\ldots$
交替排列，相邻距离为 $a$。
求链中任意一个电荷与其余电荷的相互作用能。
\exsolution
取其中一个 $+q$ 位于原点。
距离原点左右两侧各有
$n a$ 处电荷 $(-1)^nq$，$n\ge1$。
原点电荷的\emph{与其余所有电荷作用能}
为 $q$ 乘外部电势：
\[
W=2k\frac{q^2}{a}
\sum_{n=1}^\infty\frac{(-1)^n}{n}.
\]
利用
$\ln(1+x)=\sum_{n\ge1}(-1)^{n-1}x^n/n$
在 $x=1$ 的极限，
$\sum_{n\ge1}(-1)^n/n=-\ln2$，
得
\[
\boxed{W=-\frac{2q^2\ln2}
{4\pi\varepsilon_0a}}.
\]
\excheck 本题问\emph{指定一个电荷与其余电荷}
的能量。如果要求无限链的平均总能
按每个电荷分摊，应再除以2以避免成对重复计数。

\electrocapproblem{129}{109}
点电荷 $q$ 距无限接地金属导体平面为
$\ell$。求该电荷与导体表面
\emph{感应电荷的相互作用能}。
\exsolution
导体的感应电荷在真电荷位置所产生的电势
等于镜像电荷 $-q$ 在距离 $2\ell$ 处的电势：
\[
\varphi_{\rm ind}=
-\frac{q}{4\pi\varepsilon_0(2\ell)}
=-\frac{q}{8\pi\varepsilon_0\ell}.
\]
本题直接求\emph{真电荷与给定感应分布}
之间的双体相互作用项
$W_{\rm int}=q\varphi_{\rm ind}$，所以
\[
\boxed{W_{\rm int}=
-\frac{q^2}{8\pi\varepsilon_0\ell}}.
\]
\excheck 这与从无限远\emph{组装}
真电荷—导体完整体系所需的
$-\frac{q^2}{16\pi\varepsilon_0\ell}$
不相同：后者包含诱导分布随运动改变的
$1/2$ 能量因子。原书本题明确问相互作用项。

\electrocapproblem{130}{109--110}
两颗球形对称带电小球，净电荷为
$q_1,q_2$，两球中心距离为
$\ell$，彼此不重叠。
求两球间相互作用能。
\exsolution
对球外点，球对称电荷分布产生的静电势
与中心点电荷相同：
$\varphi_1(r)=kq_1/r$。
如果第二球完整位于第一球之外，
第一球的势在第二球的各球壳上按球对称平均，
等于其中心位置的势
$kq_1/\ell$（平均值定理）。
因此
\[
W_{12}=\int_{\text{球2}}\varphi_1(\bm r)
\,\mathrm dq_2(\bm r)
=q_2\frac{kq_1}{\ell},
\]
故
\[
\boxed{W_{12}=
\frac{q_1q_2}{4\pi\varepsilon_0\ell}}.
\]
\excheck 需两颗球的电荷分布均保持球对称；
若球是相互极化的导体，近距离时分布可能改变，
不能直接套此简式。
% === END 18-capacitance-energy-122-130.tex ===

% === BEGIN 18-capacitance-energy-131-140.tex blob 8bc587856a5bc21541ab3a279ce0045f1ecdf039 ===
% Irodov Part III §3.3 131--140; source PDF pp.119--120.
% The PDF printed page 110 repeats the title of problem 130 as an extra "131";
% the second "131" is the capacitor problem matched by the original answer 131.
\electrocapproblem{131}{110}
电容为 $C_1=1.0\,\mathrm{\mu F}$、
初始电压 $V=300\,\mathrm V$ 的电容器，
与一只未充电的电容器
$C_2=2.0\,\mathrm{\mu F}$ 并联相接。
求系统达到平衡后的电场能增量，并解释能量去向。
\exsolution
连接前总电荷
$Q_0=C_1V$，总电场能
$W_0=\frac12C_1V^2$。
并联后两只电容器电势相等且总电荷守恒：
\[
V_f=\frac{Q_0}{C_1+C_2}=\frac{C_1V}{C_1+C_2}.
\]
末态电场能
$W_f=\frac12(C_1+C_2)V_f^2
=\frac12 C_1^2V^2/(C_1+C_2)$，
所以
\[
\boxed{\Delta W=W_f-W_0
=-\frac12V^2\frac{C_1C_2}{C_1+C_2}}.
\]
数值：
\[
\boxed{\Delta W=-0.030\,\mathrm J=-30\,\mathrm{mJ}}.
\]
消失的电场能在实际连接导线中
变为焦耳热、电磁辐射等，
不会凭空消失。
\excheck 原版 PDF 的正文将第130题重复印成额外
“131”，此题是第二次出现的“131”且与答案131对应。
书后数值若显示“$0.03\,\mathrm{mJ}$”
则与 $1\mu\mathrm F$、300V 的计算相差1000倍；
本稿保留正确 SI 算值并标勘误。

\electrocapproblem{132}{110}
图3.30：两只电容 $C$、
中间电容 $C_0$ 与电池 $\mathcal E$
构成电路。开关 Sw 从位置1拨到2。
求开关过程中产生的焦耳热。
\exsolution
把开关前后相应节点看作
两个通过电容器耦合的等势导体，
对于固定电源电动势，
电容电路的\emph{电势跳变}激发一次电荷再分配。
图3.30左右两个相等电容 $C$
从公共上导体接到由电池保持电压差
$\mathcal E$ 的下方两节点。
设开关连接的中间电容 $C_0$
下端初连左节点、末连右节点。
切换引起的中间支路电位跃变为 $\mathcal E$；
消除其公共模式后，
参与再分配的等效电容
\[
C_{\rm tr}=\frac{2CC_0}{2C+C_0}.
\]
一次从静态到静态的切换
消耗的电场—电源可用能
$\frac12C_{\rm tr}(\Delta V)^2$，
于是
\[
\boxed{Q=\frac{CC_0}{2C+C_0}\,\mathcal E^2}.
\]
\excheck 本题求的是有微小实际电阻时最终释放的
总焦耳热，理想零电阻模型只提供初末静态电荷；
若开关瞬时切换，过程可含瞬态电磁辐射。

\electrocapproblem{133}{110}
图3.31：一只电容 $C$ 通过换向开关
先接与电动势 $\mathcal E_1$ 对应的节点，
再从位置1切至位置2，
后者改变电容两端的电位差，
求电路释放的热量。
\exsolution
设切换前电容器电压为 $U_i$，
切换后电压为 $U_f$。
对于理想电源维持的最终电压状态，
通过微弱电阻缓慢耗散的电荷再分配过程，
其热量满足普适恒等式
\[
Q=\frac12C(U_f-U_i)^2.
\]
由原书图3.31两个源的串联位置及开关端点，
开关从1切到2时电容器电压\emph{净改变}
$U_f-U_i=\mathcal E_2$，
所以
\[
\boxed{Q=\frac12C\mathcal E_2^2}.
\]
\excheck 结果与 $\mathcal E_1$ 无关，
因为焦耳热只由切换产生的电压差平方决定；
本结论使用图中电源的具体极性与接线。

\electrocapproblem{134}{110}
两层同心薄金属球壳半径
$R_1<R_2$，电荷分别为 $q_1,q_2$。
求各球壳静电自能 $W_1,W_2$，
两球壳相互作用能 $W_{12}$，
以及总电场能。
\exsolution
单个薄球壳电场能为
$W_{\rm self}=q^2/(8\pi\varepsilon_0R)$，
所以
\[
\boxed{W_1=\frac{q_1^2}{8\pi\varepsilon_0R_1}},
\qquad
\boxed{W_2=\frac{q_2^2}{8\pi\varepsilon_0R_2}}.
\]
外球壳在内球壳位置的电势恒为
$kq_2/R_2$，故交互能
\[
\boxed{W_{12}=\frac{q_1q_2}{4\pi\varepsilon_0R_2}}.
\]
相加：
\[
\boxed{W=
\frac1{8\pi\varepsilon_0}
\left(\frac{q_1^2}{R_1}+
\frac{q_2^2+2q_1q_2}{R_2}\right)}.
\]
\excheck 交互能只计入一次，
不应在 $W_1,W_2$ 内重复算交互项。

\electrocapproblem{135}{110}
半径为 $R$ 的实心绝缘球，
体电荷均匀且总电荷为 $q$，
周围真空。求总静电自能
及球内储能 $W_{\rm in}$ 与球外储能
$W_{\rm out}$ 的比值。
\exsolution
高斯定律给球内
$E(r)=qr/(4\pi\varepsilon_0R^3)$，
球外
$E(r)=q/(4\pi\varepsilon_0r^2)$。
场能密度
$w=\varepsilon_0E^2/2$。
球内积分：
\[
W_{\rm in}=\int_0^R
\frac{\varepsilon_0}{2}
\frac{q^2r^2}{16\pi^2\varepsilon_0^2R^6}
4\pi r^2\,\mathrm dr
=\frac{q^2}{40\pi\varepsilon_0R}.
\]
球外积分：
\[
W_{\rm out}=\int_R^\infty
\frac{\varepsilon_0}{2}
\frac{q^2}{16\pi^2\varepsilon_0^2r^4}
4\pi r^2\,\mathrm dr
=\frac{q^2}{8\pi\varepsilon_0R}.
\]
于是
\[
\boxed{W=W_{\rm in}+W_{\rm out}
=\frac{3q^2}{20\pi\varepsilon_0R}},
\qquad
\boxed{\frac{W_{\rm in}}{W_{\rm out}}=\frac15}.
\]
\excheck 球体内能只占总场能 $1/6$；
若把实心球按金属球壳处理将漏掉球内能。

\electrocapproblem{136}{110}
点电荷 $q=3.0\,\mathrm{\mu C}$
位于球心，
$R_1=a=250\,\mathrm{mm}$ 至
$R_2=b=500\,\mathrm{mm}$ 的球壳空间
填满相对介电常数 $\varepsilon=3.0$
的均匀介质。
求仅存储在这层\emph{介质体积}内的静电场能。
\exsolution
介质中的径向位移
$D=q/(4\pi r^2)$，
电场
$E=q/(4\pi\varepsilon_0\varepsilon r^2)$。
场能密度
$w=\frac12\bm E\cdot\bm D$，
故
\[
W=\int_a^b \frac{q^2}
{32\pi^2\varepsilon_0\varepsilon r^4}
4\pi r^2\,\mathrm dr
=\frac{q^2}{8\pi\varepsilon_0\varepsilon}
\left(\frac1a-\frac1b\right).
\]
数值
\[
\boxed{W\simeq2.7\times10^{-2}\,\mathrm J
=27\,\mathrm{mJ}}.
\]
\excheck 本题只求介质\emph{球壳内}场能，
不包括 $r<a$ 与 $r>b$ 空间的能量。

\electrocapproblem{137}{110}
均匀带电薄球壳总电荷固定为 $q$，
半径从 $R_1$ 缓慢扩大为 $R_2>R_1$。
求静电斥力做功。
\exsolution
均匀带电薄球壳自能为
$W(R)=q^2/(8\pi\varepsilon_0R)$。
自由电荷守恒，静电力做的功
$A=W(R_1)-W(R_2)$，故
\[
\boxed{A=\frac{q^2}{8\pi\varepsilon_0}
\left(\frac1{R_1}-\frac1{R_2}\right)}.
\]
\excheck 扩张时场能下降，电力做正功；
这里忽略机械表面张力和其它非电作用。

\electrocapproblem{138}{111}
总电荷固定为 $q$ 的均匀薄球壳
半径从 $R_1$ 增至 $R_2$，
球心另外放有点电荷 $q_0$。
求电场力对球壳扩张所做功。
\exsolution
壳自身能量
$W_{\rm shell}=q^2/(8\pi\varepsilon_0R)$。
球心点电荷与壳的相互作用能为
$W_{\rm int}=q_0\varphi_{\rm shell}(0)
=q_0q/(4\pi\varepsilon_0R)$。
点电荷自能独立于 $R$，在能量差中相消。
因此
\[
A=[W_{\rm shell}(R_1)+W_{\rm int}(R_1)]
-[W_{\rm shell}(R_2)+W_{\rm int}(R_2)].
\]
整理：
\[
\boxed{A=
\frac{q(q+2q_0)}{8\pi\varepsilon_0}
\left(\frac1{R_1}-\frac1{R_2}\right)}.
\]
\excheck 若 $q_0$ 与 $q$ 异号且足够大，
作用力可能反而阻碍扩张，此时 $A<0$。

\electrocapproblem{139}{111}
均匀带电薄球壳的面电荷密度为 $\sigma$。
利用能量守恒求其表面每单位面积所受
向外静电压力（法向力模）。
\exsolution
半径为 $R$ 的球壳总电荷
$q=4\pi R^2\sigma$。
以总电荷 $q$ 固定为前提，
球壳自能
$W=q^2/(8\pi\varepsilon_0R)$。
若半径增大 $\mathrm dR$，
体积增量
$\mathrm dV=4\pi R^2\mathrm dR$。
向外压力 $p_e$ 做功
$p_e\mathrm dV=-\mathrm dW$，所以
\[
p_e=-\frac1{4\pi R^2}\frac{\mathrm dW}{\mathrm dR}
=\frac{q^2}{32\pi^2\varepsilon_0R^4}.
\]
利用
$\sigma=q/(4\pi R^2)$：
\[
\boxed{p_e=\frac{\sigma^2}{2\varepsilon_0}}.
\]
\excheck 变形求导时应保持总电荷 $q$
不变，不能误保持面电荷密度 $\sigma$ 不变。

\electrocapproblem{140}{111}
球心 O 放有点电荷 $q$，
其外有一层整体不带净电荷的球形导体壳，
内外半径为 $a<b$，并有小孔（图3.32）。
求通过小孔把电荷从中心缓慢移至无穷远
需要的外功。
\exsolution
以导体壳的感应电荷解为基础：
中心点电荷使内球面感应 $-q$，
外球面感应 $+q$。
初态电场在 $0<r<a$ 和 $r>b$
分别存在，中间导体内部为零。
若将点电荷移至无穷远，
球壳总电荷仍为零，极限态无电场。
两态与点电荷自身发散自能无关的
\emph{可变有限场能}差，
可由空心球壳在两个半径之间的
缺失场能计算：
\[
A=\frac12\int_a^b
\frac{q^2}{(4\pi\varepsilon_0r^2)^2}
\varepsilon_0(4\pi r^2)\,\mathrm dr.
\]
故
\[
\boxed{A=\frac{q^2}{8\pi\varepsilon_0}
\left(\frac1a-\frac1b\right)}.
\]
\excheck 假设小孔足够小，不改变球对称场近似；
真实经过孔口时局部场会畸变，
但端点能量差在该理想极限不变。
% === END 18-capacitance-energy-131-140.tex ===

% === BEGIN 18-capacitance-energy-141-146.tex blob 3858577a02e7c085fc4bf32b83b7fd02cf63533f ===
% Irodov Part III §3.3 final problems 141--146;
% source PDF pp.120--121 (printed 111--112); figures 3.33, 3.34.
\electrocapproblem{141}{111}
空气平板电容器的极板面积均为 $S$，
缓慢把板间距从 $x_1$ 增至 $x_2>x_1$。
求所需的外界机械功：
(a) 极板自由电荷 $q$ 保持不变；
(b) 两极板电压 $V$ 保持不变。
\exsolution
忽略边缘场，
$C(x)=\varepsilon_0S/x$。
(a) 固定自由电荷时
$U(q,x)=q^2/(2C)=q^2x/(2\varepsilon_0S)$。
外界准静态功等于场能增加量：
\[
\boxed{A_q=\frac{q^2(x_2-x_1)}{2\varepsilon_0S}}.
\]
(b) 固定电压时电源会交换电荷，
不能简单取储能差作外功。
极板间吸引压力为
$p_e=\varepsilon_0V^2/(2x^2)$，
力 $F=p_eS$，
因此
\[
A_V=\int_{x_1}^{x_2}\frac{\varepsilon_0SV^2}{2x^2}
\,\mathrm dx
=\boxed{\frac{\varepsilon_0SV^2}{2}
\left(\frac1{x_1}-\frac1{x_2}\right)}.
\]
\excheck 原书(a)英文正文有
“capacitance ... equal to q”的明显笔误，
$q$ 量纲为电荷，依物理意义必须是
“charge $q$ remains constant”；
与原书答案的 $q^2$ 表达式一致。

\electrocapproblem{142}{111}
平板电容器原电容 $C_0=20\,\mathrm{nF}$，
极板间隙为 $d$。
先插入厚度为 $h=\eta d$ 的平行板
（$\eta=0.60$），接电压
$V=200\,\mathrm V$ 充电，随后切断电源，
再缓慢抽出该插板。
求抽出所需机械功：
(a) 金属板；(b) 玻璃介质板。
\exsolution
设充电时插板在位的电容为 $C_i$。
电源断开后自由电荷
$Q=C_iV$ 保持不变；
抽出后电容回到 $C_0$。
因此机械功
\[
A=\Delta U=
\frac{Q^2}{2}\left(\frac1{C_0}-\frac1{C_i}\right)
=\frac{C_i^2V^2}{2}
\left(\frac1{C_0}-\frac1{C_i}\right).
\]
(a) 金属板内部电场为零，
只剩有效气隙厚度
$d-h=d(1-\eta)$，
故 $C_i=C_0/(1-\eta)$，得到
\[
\boxed{A_{\rm metal}=
\frac12C_0V^2\frac{\eta}{(1-\eta)^2}
=1.5\,\mathrm{mJ}}.
\]
(b) 玻璃相对介电常数为 $\varepsilon$，
与两侧空气为法向串联。
设
$f=1-\eta+\eta/\varepsilon$，
则 $C_i=C_0/f$，所以
\[
\boxed{A_{\rm glass}=
\frac12C_0V^2
\frac{\eta(1-\varepsilon^{-1})}
{(1-\eta+\eta/\varepsilon)^2}}.
\]
\excheck 第(b)问题面未给玻璃 $\varepsilon$ 数值，
应保留参数表达式，
不可擅自代入某种玻璃的介电常数。

\electrocapproblem{143}{111}
将水平放置的平板电容器浸没在水中，
极板间距 $d=1.0\,\mathrm{mm}$，
水完全填满间隙。
两极板电压保持为 $V=500\,\mathrm V$。
求水在间隙中的电压力增量。
\exsolution
取水的相对介电常数为
$\varepsilon$，忽略边缘效应。
水中电场 $E=V/d$。
界面电场能和介质极化在固定外加电势的
电机械耦合关系给水中附加的电压力
\[
\Delta p=\frac12\varepsilon_0
\varepsilon(\varepsilon-1)E^2.
\]
代入 $E=V/d$：
\[
\boxed{\Delta p=
\frac{\varepsilon_0\varepsilon(\varepsilon-1)V^2}{2d^2}}.
\]
室温水取 $\varepsilon\simeq81$：
\[
\boxed{\Delta p\simeq7.2\times10^3\,\mathrm{Pa}
\approx7\,\mathrm{kPa}\approx0.07\,\mathrm{atm}}.
\]
\excheck 该式是本题所采用的水介电响应
电压与压力模型，
与只计算真空 Maxwell 电压力
$\varepsilon_0E^2/2$ 不同。

\electrocapproblem{144}{111}
水平平板电容器的一块极板浸在液体内，
另一块位于液面上方（图3.33）。
液体密度 $\rho$，相对介电常数
$\varepsilon$。
极板表面自由电荷密度模为 $\sigma$，
求液体在两极板之间升高的高度 $h$。
\exsolution
两板间浸入介质的部分增加时，
其单位体积静电储能由
$u_{\rm vac}=\sigma^2/(2\varepsilon_0)$
变为
$u_{\rm diel}=\sigma^2/(2\varepsilon_0\varepsilon)$。
在保持自由电荷密度的准静态近似下，
每单位上升体积可释放的电场能为
\[
\Delta u=\frac{\sigma^2}{2\varepsilon_0}
\left(1-\frac1\varepsilon\right).
\]
与上升液柱静水压头 $\rho gh$ 平衡，
\[
\rho gh=\Delta u.
\]
因此
\[
\boxed{h=\frac{(\varepsilon-1)\sigma^2}
{2\varepsilon_0\varepsilon\rho g}}.
\]
\excheck 如果改为固定\emph{极板电压}，
电荷密度会随介质分布改变，
必须使用相应不同的热力学约束。

\electrocapproblem{145}{111}
同轴圆柱电容器，
两电极平均半径为 $R$，
径向间隙 $d\ll R$。
向间隙沿轴向插入相对介电常数
$\varepsilon$ 的圆柱形介质层，
两电极保持恒定电压 $V$。
求把介质进一步吸入的电力。
\exsolution
忽略端部效应，插入轴向长度为 $x$
时增加的介质体积为
$2\pi Rd\,x$。
这相当于单位插入长度新增并联电容
\[
\frac{\mathrm dC}{\mathrm dx}
\simeq\frac{\varepsilon_0(\varepsilon-1)}d
(2\pi R).
\]
在固定电压下，
电源—电场系统的有效势为
$-\frac12CV^2$，
故吸入力
\[
F=\frac12V^2\frac{\mathrm dC}{\mathrm dx}.
\]
于是
\[
\boxed{F=\frac{\pi R\varepsilon_0
(\varepsilon-1)V^2}{d}}.
\]
\excheck 这是介质已部分插入时
接近轴向端缘的电力，
不表示介质无限延伸后的
横向均匀区域仍有非零吸入力。

\electrocapproblem{146}{112}
图3.34：固定电极由两片半径为 $R$
的半圆形金属极板组成，
板间距 $d$。薄可转动介质板绕
半圆中心 O 转动，
相对介电常数为 $\varepsilon$，
转角为 $\phi$，固定电压 $V$。
求绕 O 的电力矩模。
\exsolution
忽略边缘场，
介质覆盖的两极板重叠面积随转角改变。
半径 $R$ 的圆扇形面积为
$A(\phi)=\frac12R^2\phi$，
所以
$\mathrm dA/\mathrm d\phi=R^2/2$。
被介质覆盖区域相对于真空的
增量电容为
$\Delta C=(\varepsilon-1)\varepsilon_0A/d$，
因此
\[
\frac{\mathrm dC}{\mathrm d\phi}
=\frac{\varepsilon_0(\varepsilon-1)R^2}{2d}.
\]
固定电压的电力矩
$\mathcal N=\frac12V^2\mathrm dC/\mathrm d\phi$，
故
\[
\boxed{|\mathcal N|=
\frac{\varepsilon_0(\varepsilon-1)R^2V^2}{4d}}.
\]
\excheck 在原图允许的旋转角度范围内，
扇形覆盖面积随 $\phi$ 线性变化，
故电力矩大小不依赖 $\phi$。
% === END 18-capacitance-energy-141-146.tex ===
